\begin{longtable}{@{}p{3.5cm}p{2cm}p{2cm}p{2cm}p{2cm}p{2cm}p{2.5cm}@{}}
\caption{Classification des enveloppes stellaires basée sur l'analyse de résolution} \\
\toprule
\textbf{Star} & \textbf{Filter} & \textbf{FWHM\_Star} & \textbf{FWHM\_PSF} & \textbf{$\eta$} & \textbf{Has\_PSF} & \textbf{Statut de\\l'enveloppe} \\
\midrule
\endfirsthead

\multicolumn{7}{c}%
{{\bfseries \tablename\ \thetable{} -- Suite de la page précédente}} \\
\toprule
\textbf{Star} & \textbf{Filter} & \textbf{FWHM\_Star} & \textbf{FWHM\_PSF} & \textbf{$\eta$} & \textbf{Has\_PSF} & \textbf{Statut de\\l'enveloppe} \\
\midrule
\endhead

\midrule
\multicolumn{7}{r}{{Suite à la page suivante}} \\
\endfoot

\bottomrule
\multicolumn{7}{p{\textwidth}}{%
\textbf{Notes :} \\
$\eta$ : Rapport entre la FWHM de l'étoile et celle de la PSF ($\eta = \text{FWHM}_{\text{étoile}}/\text{FWHM}_{\text{PSF}}$). \\
\textbf{CR} : Clairement résolu ($\eta \geq 1.35$). \\
\textbf{MR} : Marginalement résolu ($1 < \eta < 1.35$). \\
\textbf{NR} : Non résolu ($\eta \leq 1$). \\
Has\_PSF indique si l'étoile possède sa propre PSF de référence.%
} \\
\endlastfoot

% LES DONNÉES SERONT INSÉRÉES ICI PAR LE SCRIPT PYTHON
% \input{Tables/all_stars_classification.tex}

\end{longtable}
